% ==============================================================================
% Figura: Diagrama de Flujo Lógico de la Multiplexación de Estados (1.0 kHz)
% Refleja exactamente Furuta_ControlISR_1kHz + Furuta_ProcessCommand + TIM6 ISR
% en furuta_control.c y stm32f4xx_it.c, cubriendo los 10 estados de Furuta_State_t.
% Estrategia: Layout by Construction, paleta idéntica a Figs. 5-11, cero cruces.
% Bibliotecas TikZ requeridas: arrows.meta, positioning, fit, calc,
%                              shapes.geometric, backgrounds, babel, bending.
% ==============================================================================
\begin{figure*}[!t]
\centering
% --- Paleta semántica armónica (idéntica a Figs. 5-11) ---
\definecolor{flInk}{HTML}{1E293B}%
\definecolor{flMuted}{HTML}{475569}%
\definecolor{flIdleBdr}{HTML}{243447}%
\definecolor{flIdleFill}{HTML}{ECEFF4}%
\definecolor{flServBdr}{HTML}{33658A}%
\definecolor{flServFill}{HTML}{EEF5FA}%
\definecolor{flArmBdr}{HTML}{0F766E}%
\definecolor{flArmFill}{HTML}{E6F7F5}%
\definecolor{flSwBdr}{HTML}{B45309}%
\definecolor{flSwFill}{HTML}{FEF6EB}%
\definecolor{flBalBdr}{HTML}{1B4965}%
\definecolor{flBalFill}{HTML}{EAF4FB}%
\definecolor{flActBdr}{HTML}{A63429}%
\definecolor{flActFill}{HTML}{FDF0EE}%
\definecolor{flMuxBdr}{HTML}{7E2268}%
\definecolor{flMuxFill}{HTML}{FBF0F8}%
\definecolor{flDecBdr}{HTML}{B45309}%
\definecolor{flDecFill}{HTML}{FFFBEB}%

% --- Parámetros centralizados de geometría (Single Source of Truth) ---
\def\flCaseW{2.85cm}%
\def\flRowH{0.58cm}%
\def\flSubDrop{0.84cm}%
\def\flActW{2.45cm}%

\resizebox{\textwidth}{!}{%
\begin{tikzpicture}[
    font=\sffamily,
    % --- Estilos ISO 5807 armonizados ---
    term/.style={
        rectangle,
        rounded corners=10pt,
        draw=flIdleBdr,
        fill=flIdleFill,
        line width=1.25pt,
        minimum height=0.66cm,
        inner xsep=8pt,
        inner ysep=3pt,
        align=center,
        text=flInk,
        font=\footnotesize\bfseries\sffamily
    },
    proc/.style={
        rectangle,
        rounded corners=3.5pt,
        draw=#1,
        fill=#1!8!white,
        line width=0.95pt,
        minimum height=\flRowH,
        inner xsep=5pt,
        inner ysep=2.5pt,
        align=center,
        text=flInk,
        font=\footnotesize\sffamily
    },
    casepill/.style={
        rectangle,
        rounded corners=4pt,
        draw=#1,
        fill=#1!12!white,
        line width=1.10pt,
        text width=\flCaseW,
        minimum height=\flRowH,
        inner xsep=4pt,
        inner ysep=2.5pt,
        align=center,
        text=flInk,
        font=\footnotesize\bfseries\sffamily
    },
    muxbox/.style={
        rectangle,
        rounded corners=5pt,
        draw=flMuxBdr,
        fill=flMuxFill,
        line width=1.35pt,
        text width=3.10cm,
        minimum height=0.80cm,
        inner xsep=5pt,
        inner ysep=3pt,
        align=center,
        text=flMuxBdr,
        font=\footnotesize\bfseries\sffamily
    },
    actout/.style={
        rectangle,
        rounded corners=3.5pt,
        draw=flActBdr,
        fill=flActFill,
        line width=0.95pt,
        text width=\flActW,
        minimum height=0.52cm,
        inner xsep=4pt,
        inner ysep=2pt,
        align=center,
        text=flInk,
        font=\scriptsize\bfseries\sffamily
    },
    transout/.style={
        rectangle,
        rounded corners=3.5pt,
        draw=flMuxBdr,
        fill=flMuxFill,
        line width=0.95pt,
        minimum height=0.52cm,
        inner xsep=5pt,
        inner ysep=2pt,
        align=center,
        text=flMuxBdr,
        font=\scriptsize\bfseries\sffamily
    },
    dec/.style={
        diamond,
        aspect=2.35,
        draw=flDecBdr,
        fill=flDecFill,
        line width=1.0pt,
        inner xsep=2pt,
        inner ysep=1.2pt,
        minimum width=1.95cm,
        minimum height=0.72cm,
        align=center,
        text=flInk,
        font=\scriptsize\bfseries\sffamily
    },
    zonebox/.style={
        rectangle,
        rounded corners=6pt,
        draw=#1,
        densely dashed,
        line width=0.9pt,
        fill=#1,
        fill opacity=0.035,
        text opacity=1
    },
    zonetitle/.style={
        rectangle,
        rounded corners=3pt,
        draw=#1,
        line width=0.75pt,
        fill=white,
        inner xsep=5pt,
        inner ysep=1.8pt,
        font=\scriptsize\bfseries\sffamily,
        text=#1
    },
    % --- Flechas canónicas ---
    flw/.style={-{Stealth[length=4.0pt, width=2.9pt]}, line width=0.95pt, draw=flInk},
    flwshort/.style={-{Stealth[length=3.5pt, width=2.6pt]}, line width=0.95pt, draw=flInk},
    flwline/.style={line width=0.95pt, draw=flInk},
    flwmux/.style={-{Stealth[length=4.2pt, width=3.0pt]}, line width=1.05pt, densely dashed, draw=flMuxBdr},
    flwbyp/.style={-{Stealth[length=4.2pt, width=3.0pt]}, line width=0.95pt, densely dashed, draw=flServBdr},
    dot/.style={circle, fill=#1, inner sep=1.35pt},
    lbl/.style={font=\scriptsize\sffamily, text=flInk, inner sep=1.2pt}
]

% ==============================================================================
% 1. BANDA SUPERIOR: ENTRADA ISR TIM6 (1 kHz), GUARDA CALIBRATING Y SUPERLAZO
% ==============================================================================
\node[term] (isr_start) at (0, 0) {ISR \texttt{TIM6} (1 kHz)};
\node[dec, minimum width=2.45cm, right=0.85cm of isr_start] (dec_cal) {¿\texttt{CALIBRATING}?};
\node[proc=flArmBdr, text width=2.75cm, right=1.05cm of dec_cal] (acq_sens)
    {\textbf{Lectura AS5600}\\$\theta,\dot{\theta}_{f1},\alpha,\tilde{\alpha},\dot{\alpha}_{f1}$};

% Superlazo de fondo asíncrono (UART DMA / B1 -> Parser CLI)
\node[proc=flMuxBdr, text width=2.35cm, right=1.35cm of acq_sens] (uart_b1) {\textbf{UART DMA / B1}};
\node[proc=flMuxBdr, text width=2.25cm, right=0.85cm of uart_b1] (parser) {\textbf{Parser CLI}};

\begin{scope}[on background layer]
    \node[zonebox=flMuxBdr, fit=(uart_b1)(parser), inner xsep=8pt, inner ysep=8pt] (zBg) {};
    \node[zonetitle=flMuxBdr, anchor=west] at ([xshift=0.25cm]zBg.north west) {Superlazo Fondo (as\'incrono)};
\end{scope}

% Despachador lógico central: switch(sys->state)
\node[muxbox] (mux_sw) at (0, -1.55) {MUX L\'ogico\\\texttt{switch(sys->state)}};

% Conexiones banda superior (cero cruces)
\draw[flwshort] (isr_start.east) -- node[lbl, midway, above=1.5pt]{\texttt{tick++}} (dec_cal.west);
\draw[flwshort] (dec_cal.east) -- node[lbl, midway, above=1.5pt]{\textbf{No}} (acq_sens.west);
\draw[flwshort, draw=flMuxBdr] (uart_b1.east) -- (parser.west);

% Lectura AS5600 -> MUX (carril superior de entrada a mux_sw.east)
\draw[flw, rounded corners=4pt]
    (acq_sens.south) |- ([yshift=0.16cm]mux_sw.east);

% Parser CLI -> MUX (carril inferior punteado, estrictamente debajo del de AS5600)
\draw[flwmux, rounded corners=4pt]
    (parser.south) |- node[lbl, pos=0.66, above=2pt, text=flMuxBdr]{\textbf{\texttt{sel / reset}} (actualiza \texttt{sys->state})}
    ([yshift=-0.18cm]mux_sw.east);

% ==============================================================================
% 2. REJILLA DE COORDENADAS X PARA LAS 8 RAMAS DEL MULTIPLEXOR (SINGLE SOURCE)
% ==============================================================================
\coordinate (x_spine)  at (0, 0);
\coordinate (x_case)   at (0.72, 0);
\coordinate (x_c1)     at (4.35, 0);
\coordinate (x_d1)     at (5.45, 0);
\coordinate (x_d2)     at (8.65, 0);
\coordinate (x_d3)     at (11.80, 0);
\coordinate (x_out)    at (18.15, 0);
\coordinate (x_coll)   at (20.25, 0);
\coordinate (x_bypass) at (20.98, 0);

% Coordenadas Y de las 8 filas (ritmo determinista sin solapamientos)
\def\yRone{-3.15}
\def\yRtwo{-4.25}
\def\yRthree{-5.70}
\def\yRfour{-6.80}
\def\yRfive{-8.30}
\def\yRsix{-9.80}
\def\yRseven{-11.50}
\def\yReight{-14.05}

% ==============================================================================
% 3. FILAS 1-6: ESTADOS PASIVOS, SERVICIO, MOTOR AISLADO, PARADA Y RECUPERACIÓN
% ==============================================================================

% --- FILA 1: IDLE / FAULT_STOP (Estados 0 y 7) ---
\node[casepill=flIdleBdr, anchor=west] (c1) at (x_case |- 0,\yRone) {IDLE / FAULT\_STOP};
\node[proc=flIdleBdr, text width=2.85cm, anchor=west] (p1) at (x_c1 |- c1) {$u = 0$, \texttt{slew} $= 0$};
\node[actout, anchor=center] (o1) at (x_out |- c1) {\texttt{L298N\_Coast}};
\draw[flwshort] (c1.east) -- (p1.west);
\draw[flw] (p1.east) -- (o1.west);
\draw[flwline] (o1.east) -- (x_coll |- o1);

% --- FILA 2: CHARACTERIZATION (Estado 2) ---
\node[casepill=flServBdr, anchor=west] (c2) at (x_case |- 0,\yRtwo) {CHARACTERIZATION};
\node[proc=flServBdr, text width=2.85cm, anchor=west] (p2) at (x_c1 |- c2) {Ensayo \texttt{CHAR1--5B}};
\node[dec, anchor=center] (d2) at (x_d2 |- c2) {¿Fin?};
\node[transout, text width=2.65cm, anchor=center] (o2a) at (x_out |- d2) {\texttt{state} $\leftarrow$ \texttt{IDLE}};
\node[actout, anchor=center] (o2b) at ($(x_out |- d2)+(0,-\flSubDrop)$) {PWM / Coast};
\draw[flwshort] (c2.east) -- (p2.west);
\draw[flwshort] (p2.east) -- (d2.west);
\draw[flw] (d2.east) -- node[lbl, pos=0.14, above=1.5pt]{\textbf{S\'i}} (o2a.west);
\draw[flw, rounded corners=4pt] (d2.south) |- node[lbl, pos=0.22, right=2.5pt]{\textbf{No}} (o2b.west);
\draw[flwline] (o2a.east) -- (x_coll |- o2a);
\draw[flwline] (o2b.east) -- (x_coll |- o2b);

% --- FILA 3: MOTOR_PID_VEL (Estado 3) ---
\node[casepill=flArmBdr, anchor=west] (c3) at (x_case |- 0,\yRthree) {MOTOR\_PID\_VEL};
\node[proc=flArmBdr, text width=2.85cm, anchor=west] (p3) at (x_c1 |- c3) {PID Velocidad};
\node[actout, anchor=center] (o3) at (x_out |- c3) {\texttt{L298N\_SetVoltage}};
\draw[flwshort] (c3.east) -- (p3.west);
\draw[flw] (p3.east) -- (o3.west);
\draw[flwline] (o3.east) -- (x_coll |- o3);

% --- FILA 4: MOTOR_PID_POS (Estado 4) ---
\node[casepill=flArmBdr, anchor=west] (c4) at (x_case |- 0,\yRfour) {MOTOR\_PID\_POS};
\node[proc=flArmBdr, text width=2.85cm, anchor=west] (p4) at (x_c1 |- c4) {Servo Clegg ($\theta_{\mathrm{ref}}$)};
\node[dec, anchor=center] (d4) at (x_d2 |- c4) {¿Bloq.?};
\node[actout, anchor=center] (o4a) at (x_out |- d4) {\texttt{L298N\_Brake}};
\node[actout, anchor=center] (o4b) at ($(x_out |- d4)+(0,-\flSubDrop)$) {\texttt{L298N\_SetVoltage}};
\draw[flwshort] (c4.east) -- (p4.west);
\draw[flwshort] (p4.east) -- (d4.west);
\draw[flw] (d4.east) -- node[lbl, pos=0.14, above=1.5pt]{\textbf{S\'i}} (o4a.west);
\draw[flw, rounded corners=4pt] (d4.south) |- node[lbl, pos=0.22, right=2.5pt]{\textbf{No}} (o4b.west);
\draw[flwline] (o4a.east) -- (x_coll |- o4a);
\draw[flwline] (o4b.east) -- (x_coll |- o4b);

% --- FILA 5: STOP_HOME (Estado 9) ---
\node[casepill=flActBdr, anchor=west] (c5) at (x_case |- 0,\yRfive) {STOP\_HOME};
\node[proc=flActBdr, text width=2.85cm, anchor=west] (p5) at (x_c1 |- c5) {Servo $0^\circ$ (Clegg)};
\node[dec, anchor=center] (d5) at (x_d2 |- c5) {¿En $0^\circ$?};
\node[transout, text width=2.65cm, anchor=center] (o5a) at (x_out |- d5) {\texttt{IDLE} + \texttt{Coast}};
\node[actout, anchor=center] (o5b) at ($(x_out |- d5)+(0,-\flSubDrop)$) {SetVoltage / Brake};
\draw[flwshort] (c5.east) -- (p5.west);
\draw[flwshort] (p5.east) -- (d5.west);
\draw[flw] (d5.east) -- node[lbl, pos=0.14, above=1.5pt]{\textbf{S\'i}} (o5a.west);
\draw[flw, rounded corners=4pt] (d5.south) |- node[lbl, pos=0.22, right=2.5pt]{\textbf{No}} (o5b.west);
\draw[flwline] (o5a.east) -- (x_coll |- o5a);
\draw[flwline] (o5b.east) -- (x_coll |- o5b);

% --- FILA 6: RECOVER_HOME (Estado 8) ---
\node[casepill=flActBdr, anchor=west] (c6) at (x_case |- 0,\yRsix) {RECOVER\_HOME};
\node[proc=flActBdr, text width=2.85cm, anchor=west] (p6) at (x_c1 |- c6) {Subfases \texttt{F0$\to$F1$\to$F2}};
\node[dec, anchor=center] (d6) at (x_d2 |- c6) {¿Calma?};
\node[transout, text width=2.65cm, anchor=center] (o6a) at (x_out |- d6) {Reset $\to$ \texttt{SWING\_UP}};
\node[actout, anchor=center] (o6b) at ($(x_out |- d6)+(0,-\flSubDrop)$) {Coast / PWM / Brake};
\draw[flwshort] (c6.east) -- (p6.west);
\draw[flwshort] (p6.east) -- (d6.west);
\draw[flw] (d6.east) -- node[lbl, pos=0.14, above=1.5pt]{\textbf{S\'i}} (o6a.west);
\draw[flw, rounded corners=4pt] (d6.south) |- node[lbl, pos=0.22, right=2.5pt]{\textbf{No}} (o6b.west);
\draw[flwline] (o6a.east) -- (x_coll |- o6a);
\draw[flwline] (o6b.east) -- (x_coll |- o6b);

% ==============================================================================
% 4. FILAS 7-8: CICLO AUTÓNOMO CON GUARDAS I2C2, CAPTURA, TIMEOUT Y CAÍDA
% ==============================================================================

% --- FILA 7: SWING_UP (Estado 5) ---
\node[casepill=flSwBdr, anchor=west] (c7) at (x_case |- 0,\yRseven) {SWING\_UP};
\node[dec, anchor=center] (d7_i2c) at (x_d1 |- c7) {¿\texttt{I2C2} OK?};
\node[dec, anchor=center] (d7_cap) at (x_d2 |- c7) {¿Captura?};
\node[dec, anchor=center] (d7_tmo) at (x_d3 |- c7) {¿Timeout?};
\node[proc=flSwBdr, text width=2.05cm, anchor=center] (p7) at (14.95, \yRseven) {Swing-Up PD};
\node[actout, anchor=center] (o7) at (x_out |- c7) {\texttt{L298N\_SetVoltage}};

% Salidas cortocircuitadas de Row 7 (sub-fila limpia debajo de cada rombo)
\coordinate (y7_sub)  at (0, \yRseven - 1.12);
\coordinate (y7_rail) at (0, \yRseven - 1.72);
\node[actout, text width=2.10cm, anchor=center] (o7_i2c) at (d7_i2c |- y7_sub) {\texttt{L298N\_Coast}};
\node[transout, text width=2.75cm, anchor=center] (o7_cap) at (d7_cap |- y7_sub) {Bumpless $\to$ \texttt{BALANCE}};
\node[transout, text width=2.45cm, anchor=center] (o7_tmo) at (d7_tmo |- y7_sub) {\texttt{RECOVER} + \texttt{Coast}};

\draw[flwshort] (c7.east) -- (d7_i2c.west);
\draw[flwshort] (d7_i2c.east) -- node[lbl, midway, above=1.5pt]{\textbf{S\'i}} (d7_cap.west);
\draw[flwshort] (d7_cap.east) -- node[lbl, midway, above=1.5pt]{\textbf{No}} (d7_tmo.west);
\draw[flwshort] (d7_tmo.east) -- node[lbl, midway, above=1.5pt]{\textbf{No}} (p7.west);
\draw[flwshort] (p7.east) -- (o7.west);
\draw[flwline] (o7.east) -- (x_coll |- o7);

\draw[flwshort] (d7_i2c.south) -- node[lbl, midway, right=2pt]{\textbf{No}} (o7_i2c.north);
\draw[flwshort] (d7_cap.south) -- node[lbl, midway, right=2pt]{\textbf{S\'i}} (o7_cap.north);
\draw[flwshort] (d7_tmo.south) -- node[lbl, midway, right=2pt]{\textbf{S\'i}} (o7_tmo.north);

% Riel colector dedicado de Fila 7 -> bus colector derecho
\draw[flwline, rounded corners=3pt] (o7_i2c.south) |- (x_coll |- y7_rail);
\draw[flwline] (o7_cap.south) -- (o7_cap |- y7_rail);
\draw[flwline] (o7_tmo.south) -- (o7_tmo |- y7_rail);
\node[dot=flInk] at (o7_cap |- y7_rail) {};
\node[dot=flInk] at (o7_tmo |- y7_rail) {};

% --- FILA 8: BALANCE_PID (Estado 6) ---
\node[casepill=flBalBdr, line width=1.45pt, anchor=west] (c8) at (x_case |- 0,\yReight) {BALANCE\_PID};
\node[dec, anchor=center] (d8_i2c) at (x_d1 |- c8) {¿\texttt{I2C2} OK?};
\node[dec, anchor=center] (d8_fal) at (x_d2 |- c8) {¿Ca\'ida?};
\node[proc=flBalBdr, text width=4.85cm, anchor=center] (p8) at (13.65, \yReight) {\textbf{LPF 2} + \textbf{M4 Dual} + \textbf{Dual-Notch}};
\node[actout, anchor=center] (o8) at (x_out |- c8) {\texttt{L298N\_SetVoltage}};

% Salidas cortocircuitadas de Row 8
\coordinate (y8_sub)  at (0, \yReight - 1.12);
\coordinate (y8_rail) at (0, \yReight - 1.72);
\node[actout, text width=2.10cm, anchor=center] (o8_i2c) at (d8_i2c |- y8_sub) {\texttt{L298N\_Coast}};
\node[transout, text width=5.90cm, anchor=center] (o8_fal) at (11.15, \yReight - 1.12)
    {Reset $\to$ (\texttt{auto}? \texttt{RECOVER\_HOME} : \texttt{IDLE}) + \texttt{Coast}};

\draw[flwshort] (c8.east) -- (d8_i2c.west);
\draw[flwshort] (d8_i2c.east) -- node[lbl, midway, above=1.5pt]{\textbf{S\'i}} (d8_fal.west);
\draw[flwshort] (d8_fal.east) -- node[lbl, midway, above=1.5pt]{\textbf{No}} (p8.west);
\draw[flwshort] (p8.east) -- (o8.west);
\draw[flwline] (o8.east) -- (x_coll |- o8);

\draw[flwshort] (d8_i2c.south) -- node[lbl, midway, right=2pt]{\textbf{No}} (o8_i2c.north);
\draw[flwshort] (d8_fal.south) -- node[lbl, midway, right=2pt]{\textbf{S\'i}} (d8_fal |- o8_fal.north);

% Riel colector dedicado de Fila 8 -> bus colector derecho
\draw[flwline, rounded corners=3pt] (o8_i2c.south) |- (x_coll |- y8_rail);
\draw[flwline] (o8_fal.south) -- (o8_fal |- y8_rail);
\node[dot=flInk] at (o8_fal |- y8_rail) {};

% ==============================================================================
% 5. ESPINA VERTICAL DEL MULTIPLEXOR (IZQUIERDA) Y CONTENEDOR DE RAMAS
% ==============================================================================
\draw[flwline] (mux_sw.south) -- (x_spine |- c8.west);
\foreach \n in {c1, c2, c3, c4, c5, c6, c7, c8} {
    \draw[flwshort] (x_spine |- \n.west) -- (\n.west);
}
\foreach \n in {c1, c2, c3, c4, c5, c6, c7} {
    \node[dot=flInk] at (x_spine |- \n.west) {};
}

\begin{scope}[on background layer]
    \node[zonebox=flBalBdr, inner xsep=11pt, inner ysep=10pt,
          fit=(x_spine |- c1)(c1)(o1)(c8)(o8)(y8_rail)(o7_i2c)] (zMux) {};
    \node[zonetitle=flBalBdr, anchor=east]
          at ([xshift=-0.25cm]zMux.north east) {Multiplexaci\'on Determinista por Estado (\texttt{Furuta\_ControlISR\_1kHz})};
\end{scope}

% ==============================================================================
% 6. BUS COLECTOR DERECHO, BYPASS CALIBRATING, DECIMACIÓN Y SALIDA ISR
% ==============================================================================
\node[proc=flServBdr, text width=5.25cm, anchor=east] (decim) at (19.25, \yReight - 2.65)
    {\textbf{Decimaci\'on Telemetr\'ia} (\texttt{tick \% 10 == 0} $\to$ 100 Hz)};
\node[term, left=1.10cm of decim] (isr_end) {Fin ISR \texttt{TIM6}};

% Bus colector derecho (une las 8 filas hacia decim.east superior)
\coordinate (dec_in_top) at ([yshift= 0.14cm]decim.east);
\coordinate (dec_in_bot) at ([yshift=-0.14cm]decim.east);

\draw[flw, rounded corners=5pt]
    (x_coll |- o1) -- (x_coll |- dec_in_top) -- (dec_in_top);
\foreach \p in {o2a, o2b, o3, o4a, o4b, o5a, o5b, o6a, o6b, o7, o8} {
    \node[dot=flInk] at (x_coll |- \p) {};
}
\node[dot=flInk] at (x_coll |- y7_rail) {};
\node[dot=flInk] at (x_coll |- y8_rail) {};

% Desvío superior CALIBRATING (Sí -> pausa I2C, corre estrictamente por fuera de x_coll)
\coordinate (y_byp_top) at ([yshift=0.58cm]zBg.north);
\draw[flwbyp, rounded corners=6pt]
    (dec_cal.north) |- node[lbl, pos=0.62, above=2pt, text=flServBdr]{\textbf{S\'i} (pausa bus $\text{I}^2\text{C}$ durante calibraci\'on)}
    (x_bypass |- y_byp_top) -- (x_bypass |- dec_in_bot) -- (dec_in_bot);

% Salida final hacia Fin ISR TIM6
\draw[flw] (decim.west) -- (isr_end.east);

\end{tikzpicture}%
}%
\caption{Diagrama de flujo l\'ogico de la multiplexaci\'on de estados ejecutada a \qty{1.0}{\kilo\hertz} dentro de \texttt{Furuta\_ControlISR\_1kHz} (\texttt{furuta\_control.c} y \texttt{stm32f4xx\_it.c}). Cada periodo de \texttt{TIM6} verifica la guarda previa \texttt{CALIBRATING} (evitando colisiones $\text{I}^2\text{C}$ durante el promediado de cero), adquiere ambos sensores AS5600 y despacha mediante \texttt{switch(sys->state)} ---actualizado as\'incronamente por el superlazo \texttt{UART DMA / B1} $\to$ \texttt{Parser CLI}--- una \'unica rama mutuamente excluyente entre los diez estados de \texttt{Furuta\_State\_t}, aplicando protecci\'on por racha de fallos $\text{I}^2\text{C}_2$, transferencias \textit{bumpless}, conmutaci\'on autom\'atica/manual ante ca\'ida y salida directa a \texttt{L298N\_SetVoltage}, \texttt{L298N\_Brake} o \texttt{L298N\_Coast} antes de la decimaci\'on de telemetr\'ia a \qty{100}{\hertz}.}
\label{fig:flujo_multiplexacion}
\end{figure*}
